\chapter{排序算法}

\begin{introduction}
  \item 为什么要排序
  \item 冒泡排序：两两比较，大的往后走
  \item 选择排序：每次挑最小的
  \item 插入排序：像整理扑克牌
  \item 三种排序怎么选
\end{introduction}

\chaptermotto{「把乱七八糟的数据排好队，后面的事情就简单了。」}

考试成绩要排名，找最大值要先比较，查资料要按顺序找。生活里到处都需要“排好队”。

这一章只学三种最基础、最好懂的排序：冒泡排序、选择排序、插入排序。我们不讲复杂的公式，只用手把过程走一遍，再把过程翻译成 C 语言。

\section{为什么要排序}

先看一组没有排序的分数：

\begin{center}
\begin{tabular}{@{}cccccc@{}}
\toprule
90 & 65 & 78 & 100 & 82 & 55 \\
\bottomrule
\end{tabular}
\end{center}

如果老师要问“最高分是谁”，你得从头看一遍。如果数据有 1000 个，就更麻烦了。

排好序以后：

\begin{center}
\begin{tabular}{@{}cccccc@{}}
\toprule
55 & 65 & 78 & 82 & 90 & 100 \\
\bottomrule
\end{tabular}
\end{center}

最高分在最后，最低分在最前面，找起来一眼就能看到。后面学查找、统计、排名，也会方便很多。

\textbf{排序}就是把一组数据按照从小到大（或从大到小）的顺序重新排列。

\section{冒泡排序}

\subsection{先用手做一遍}

冒泡排序的规则很简单：\textbf{从左到右，两两比较，前面的比后面大就交换。}

先拿 5 个数试一试：

\begin{center}
\begin{tabular}{@{}ccccc@{}}
\toprule
5 & 2 & 9 & 1 & 7 \\
\bottomrule
\end{tabular}
\end{center}

第一轮：

\begin{itemize}
  \item 比较 5 和 2：5 比 2 大，交换，变成 \texttt{2 5 9 1 7}；
  \item 比较 5 和 9：5 比 9 小，不换，还是 \texttt{2 5 9 1 7}；
  \item 比较 9 和 1：9 比 1 大，交换，变成 \texttt{2 5 1 9 7}；
  \item 比较 9 和 7：9 比 7 大，交换，变成 \texttt{2 5 1 7 9}。
\end{itemize}

第一轮结束，最大的 9 已经“冒”到了最后面。就像水里的气泡，大的会往上冒。

第二轮再从左边开始，但最后一个 9 已经排好，不用再比。这样一轮一轮下去，直到所有数都排好。

\subsection{写成 C 语言}

\begin{codebox}{bubble.c —— 冒泡排序}
#include <stdio.h>

int main(void)
{
    int a[8] = {5, 2, 9, 1, 7, 3, 8, 6};
    int n = 8;
    int i, j, temp;

    /* 排序前 */
    for (i = 0; i < n; i++) {
        printf("%d ", a[i]);
    }
    printf("\n");

    /* 冒泡排序 */
    for (i = 0; i < n - 1; i++) {
        for (j = 0; j < n - 1 - i; j++) {
            if (a[j] > a[j + 1]) {
                temp = a[j];
                a[j] = a[j + 1];
                a[j + 1] = temp;
            }
        }
    }

    /* 排序后 */
    for (i = 0; i < n; i++) {
        printf("%d ", a[i]);
    }
    printf("\n");

    return 0;
}
\end{codebox}

\begin{outbox}{运行结果}
5 2 9 1 7 3 8 6
1 2 3 5 6 7 8 9
\end{outbox}

\subsection{这段代码在做什么}

\begin{itemize}
  \item 外层循环 \texttt{i} 表示“第几轮”；
  \item 内层循环 \texttt{j} 表示“这一轮比较到第几个”；
  \item \texttt{n - 1 - i} 的意思是：每轮结束后，后面已经有 \texttt{i} 个数排好了，不用再比较；
  \item 如果前面的数比后面的大，就借助 \texttt{temp} 交换它们。
\end{itemize}

\begin{tipbox}[交换两个数]
交换要用临时变量，就像交换两杯水需要第三个空杯子：
\begin{codebox}{swap-three.c}
temp = a[j];
a[j] = a[j + 1];
a[j + 1] = temp;
\end{codebox}
\end{tipbox}

\begin{note}
冒泡排序每轮都把当前最大的数送到后面，所以叫“冒泡”。数据少的时候很好用，数据多了会比较慢。初学阶段先把它写熟。
\end{note}

\section{选择排序}

\subsection{先用手做一遍}

选择排序的规则是：\textbf{每一轮找出最小的数，放到最前面。}

还是这 5 个数：

\begin{center}
\begin{tabular}{@{}ccccc@{}}
\toprule
5 & 2 & 9 & 1 & 7 \\
\bottomrule
\end{tabular}
\end{center}

\begin{enumerate}
  \item 从全部数里找最小的：是 1；
  \item 把 1 和第一个数 5 交换：\texttt{1 2 9 5 7}；
  \item 从剩下的 \texttt{2 9 5 7} 里找最小的：是 2，已经在第二位，不用换；
  \item 从剩下的 \texttt{9 5 7} 里找最小的：是 5，和 9 交换：\texttt{1 2 5 9 7}；
  \item 从剩下的 \texttt{9 7} 里找最小的：是 7，和 9 交换：\texttt{1 2 5 7 9}。
\end{enumerate}

排好了。

\subsection{写成 C 语言}

\begin{codebox}{selection.c —— 选择排序}
#include <stdio.h>

int main(void)
{
    int a[8] = {5, 2, 9, 1, 7, 3, 8, 6};
    int n = 8;
    int i, j, min_index, temp;

    for (i = 0; i < n; i++) {
        printf("%d ", a[i]);
    }
    printf("\n");

    for (i = 0; i < n - 1; i++) {
        min_index = i;
        for (j = i + 1; j < n; j++) {
            if (a[j] < a[min_index]) {
                min_index = j;
            }
        }
        if (min_index != i) {
            temp = a[i];
            a[i] = a[min_index];
            a[min_index] = temp;
        }
    }

    for (i = 0; i < n; i++) {
        printf("%d ", a[i]);
    }
    printf("\n");

    return 0;
}
\end{codebox}

\begin{outbox}{运行结果}
5 2 9 1 7 3 8 6
1 2 3 5 6 7 8 9
\end{outbox}

\subsection{这段代码在做什么}

\begin{itemize}
  \item \texttt{i} 表示“现在要填第几个位置”；
  \item 先假设第 \texttt{i} 个位置上的数最小，记下它的下标 \texttt{min\_index}；
  \item 用 \texttt{j} 从后面找一圈，发现更小的就更新 \texttt{min\_index}；
  \item 找完以后，如果最小数不在第 \texttt{i} 个位置，就交换过去。
\end{itemize}

选择排序的好处是\textbf{交换次数少}，每一轮最多交换一次。

\section{插入排序}

\subsection{先用手做一遍}

插入排序就像整理扑克牌：

\begin{itemize}
  \item 左手拿已经排好的牌；
  \item 右手每次摸一张新牌；
  \item 从右往左比，把新牌插到合适的位置。
\end{itemize}

例如 \texttt{5 2 9 1 7}：

\begin{enumerate}
  \item 先看第一个数 5，它自己就是有序的；
  \item 摸到 2，比 5 小，插到 5 前面：\texttt{2 5}；
  \item 摸到 9，比 5 大，放在后面：\texttt{2 5 9}；
  \item 摸到 1，比 9、5、2 都小，插到最前面：\texttt{1 2 5 9}；
  \item 摸到 7，插到 5 和 9 之间：\texttt{1 2 5 7 9}。
\end{enumerate}

\subsection{写成 C 语言}

\begin{codebox}{insertion.c —— 插入排序}
#include <stdio.h>

int main(void)
{
    int a[8] = {5, 2, 9, 1, 7, 3, 8, 6};
    int n = 8;
    int i, j, key;

    for (i = 0; i < n; i++) {
        printf("%d ", a[i]);
    }
    printf("\n");

    for (i = 1; i < n; i++) {
        key = a[i];             /* 摸到的新牌 */
        j = i - 1;

        /* 从右往左找位置，比 key 大的都往后挪 */
        while (j >= 0 && a[j] > key) {
            a[j + 1] = a[j];
            j--;
        }
        a[j + 1] = key;         /* 把新牌放进去 */
    }

    for (i = 0; i < n; i++) {
        printf("%d ", a[i]);
    }
    printf("\n");

    return 0;
}
\end{codebox}

\begin{outbox}{运行结果}
5 2 9 1 7 3 8 6
1 2 3 5 6 7 8 9
\end{outbox}

\subsection{这段代码在做什么}

\begin{itemize}
  \item 从第二个数开始，把当前数当作“新牌” \texttt{key}；
  \item 从它前面一个数开始，从右往左看；
  \item 如果前面的数比 \texttt{key} 大，就把前面的数往后挪一格；
  \item 挪到合适的位置后，把 \texttt{key} 放进去。
\end{itemize}

插入排序在数据\textbf{原本就比较有序}的时候特别快，因为不需要挪很多次。

\section{三种排序比一比}

\begin{center}
\begin{tabular}{@{}p{0.16\textwidth}p{0.24\textwidth}p{0.24\textwidth}p{0.24\textwidth}@{}}
\toprule
排序方法 & 怎么做的 & 优点 & 适合的情况 \\
\midrule
冒泡排序 & 相邻两个比较，大的往后换 & 最好理解 & 数据很少，练手 \\
选择排序 & 每轮找最小的放到前面 & 交换次数少 & 交换比较麻烦的时候 \\
插入排序 & 像抓牌，插入已排好的部分 & 数据基本有序时很快 & 数据不多、接近有序 \\
\bottomrule
\end{tabular}
\end{center}

\begin{tipbox}[初学怎么选]
初学阶段，先把\textbf{冒泡排序}写熟。它最直观，出错也容易查。等以后数据变多，再学更快的排序方法。
\end{tipbox}

\begin{note}
这三种排序在数据很多时都比较慢。真正的大数据排序，计算机会用更聪明的算法。那是后面进阶的内容，现在不用着急。
\end{note}

\section{小应用：给成绩排序}

下面把 5 个成绩从高到低排序，用的是冒泡排序。注意比较符号从 \texttt{>} 变成了 \texttt{<}：

\begin{codebox}{sort-scores.c —— 成绩从高到低}
#include <stdio.h>

int main(void)
{
    int score[5] = {78, 92, 65, 88, 70};
    int n = 5;
    int i, j, temp;

    for (i = 0; i < n - 1; i++) {
        for (j = 0; j < n - 1 - i; j++) {
            if (score[j] < score[j + 1]) {
                temp = score[j];
                score[j] = score[j + 1];
                score[j + 1] = temp;
            }
        }
    }

    printf("从高到低：");
    for (i = 0; i < n; i++) {
        printf("%d ", score[i]);
    }
    printf("\n");

    return 0;
}
\end{codebox}

\begin{outbox}{运行结果}
从高到低：92 88 78 70 65
\end{outbox}

只要把比较的“大于”改成“小于”，就能从升序变成降序。

\section{本章小结}

\begin{center}
\begin{tabular}{@{}p{0.24\textwidth}p{0.64\textwidth}@{}}
\toprule
主题 & 要点 \\
\midrule
排序 & 把数据按从小到大或从大到小的顺序排好 \\
冒泡排序 & 相邻比较，大的往后换；每轮把最大值送到末尾 \\
选择排序 & 每轮找最小值，放到前面 \\
插入排序 & 把当前数插入左边已排好的部分 \\
交换 & 需要临时变量 \texttt{temp} \\
升序/降序 & 改一下比较符号即可 \\
\bottomrule
\end{tabular}
\end{center}

\begin{tipbox}[下一章预告]
最后我们学宏定义：给一段代码起个名字，让写程序更方便。学完这一章，这本书的正文就结束啦。
\end{tipbox}

\begin{problemset}
  \item 用冒泡排序把 \texttt{5 3 8 1 2} 从小到大排序，写出每一轮的结果。
  \item 用选择排序把 \texttt{9 4 7 2 6} 从小到大排序，写出每一轮的结果。
  \item 用插入排序把 \texttt{6 2 8 4 1} 从小到大排序，写出每一步。
  \item 把本章的冒泡排序程序改成从大到小排序。
  \item 输入 10 个整数，用冒泡排序从小到大输出。
  \item 输入 10 个整数，用选择排序从大到小输出。
  \item 输入 10 个学生的成绩，用任意一种排序从高到低输出。
  \item 想一想：如果数据本来就已经排好序了，哪种排序最省事？为什么？
  \item 思考题：冒泡排序为什么要比较 \texttt{n - 1 - i} 次，而不是每次都比较 \texttt{n - 1} 次？
\end{problemset}
